% Generated draft: requires mathematical and source review.
\documentclass[11pt]{article}
\usepackage[margin=25mm]{geometry}
\usepackage{fontspec}
\setmainfont{Latin Modern Roman}
\usepackage{amsmath,amssymb,mathtools,unicode-math}
\setmathfont{Latin Modern Math}
\usepackage{enumitem,needspace,xurl,hyperref}
\setlength{\parindent}{0pt}
\setlength{\parskip}{5pt}
\setlength{\emergencystretch}{3em}
\newcommand{\N}{\mathbb{N}}
\newcommand{\Z}{\mathbb{Z}}
\newcommand{\Q}{\mathbb{Q}}
\newcommand{\R}{\mathbb{R}}
\newcommand{\C}{\mathbb{C}}
\newcommand{\F}{\mathbb{F}}
\newcommand{\A}{\mathbb{A}}
\newcommand{\PP}{\mathbb{P}}
\newcommand{\E}{\mathbb{E}}
\newcommand{\eps}{\varepsilon}
\newcommand{\dd}{\,\mathrm{d}}
\begin{document}
\section*{UnsolvedMath ID: 800973}
\textbf{Source problem code:} catalogue entry
\subsection{Definitions and mathematical statement}
A closed manifold means a compact manifold without boundary. In the standard smooth interpretation, let \(M\) be a smooth closed manifold and let \(g\) be a smooth pseudo-Riemannian metric: a symmetric, nondegenerate covariant \(2\)-tensor. A prescribed signature \((p,q)\) means that \(g\) has \(p\) positive and \(q\) negative directions at every point, where \(p+q=\dim M\). The signature is indefinite when \(p,q>0\).

A geodesic is understood to be affinely parametrized: it satisfies \(\nabla^g_{\dot\gamma}\dot\gamma=0\), where \(\nabla^g\) is the Levi-Civita connection. The metric \(g\) is geodesically complete if, for every nonempty open interval \(I\subseteq\R\) and every geodesic \(\gamma:I\to M\), there is a geodesic \(\widetilde\gamma:\R\to M\) with \(\widetilde\gamma|_I=\gamma\).

The supplied source recalls that every Riemannian metric on a closed manifold is geodesically complete, whereas for every indefinite signature there exist metrics of that signature on closed manifolds that are not geodesically complete.

\textbf{Question (Burns--Matveev, Question 54).} What geometric assumptions on a metric \(g\) on a closed manifold \(M\) imply that \(g\) is geodesically complete? The question may also be considered within a prescribed signature \((p,q)\): identify geometric hypotheses under which every closed pseudo-Riemannian manifold of that signature is geodesically complete.

This is an open-ended request for sufficient geometric hypotheses. The supplied record does not prescribe a particular hypothesis, dimension, or form of classification, and it does not explicitly specify metric regularity; the smooth setting above makes the conventional interpretation explicit.
\subsection{Short English statement}
Find geometric conditions guaranteeing that every affinely parametrized geodesic of a metric on a compact manifold without boundary extends to all real times, possibly restricting attention to one prescribed signature. The supplied background contrasts completeness in the Riemannian case with incomplete examples in every indefinite signature.
\subsection{Sources}
\begin{itemize}
\item Local UnsolvedMath registry: supplied record with permanent ID \(800973\), source code \texttt{catalogue entry}, titled \emph{Open Problems and Questions about Geodesics --- Question 54}. The registry attributes the question to Burns and Matveev and gives the proposed year as 2019.
\item Burns and Matveev, \emph{Open Problems and Questions about Geodesics} (2019, as identified in the supplied record), Question 54. Supplied source URL: \url{https://arxiv.org/abs/1308.5417}. The background records extraction by \texttt{arxiv\_latex} and an access date of 2026-10-16. These are supplied provenance details; no independent source visit or verification is claimed.
\item The registry labels the problem \texttt{open}. This is a supplied status label, not an independently verified assessment of the literature.
\end{itemize}
\end{document}
